% \documentclass{} precedes the preamble and is typically the first command in any .tex document. Every .tex document should include this command as it defines what kind of document you intend on creating. Modifiers in square brackets [] can be added in between the text ''documentclass'' and the curly brackets {} to modify font size, templates, etc. 

\documentclass[12pt,journal,compsoc]{IEEEtran}

%-----PACKAGES-------------------------------------------------------------------------------------

% Packages include extra commands that allow for additional formatting, ranging from Graphics, Math, to Alignment. The command to include packages will always look similar to \usepackage{} where the package name is within the curly brackets {}. Packages are defined in the preamble, i.e., between the \documentclass{} and \begin{document} commands.

\usepackage{graphicx}
\usepackage{textcomp}
\usepackage{graphicx}
\usepackage{amsmath}
\usepackage{hyperref}
\usepackage{booktabs}
\usepackage{geometry}

%----- The DOCUMENT Environment-------------------------------------------------------------------

% The \begin{document} and \end{document} commands establish the environment for the text of the document. The \begin{} and \end{} commands are used repeatedly in LaTeX to show where an environment begins and ends. The \end{document} command will be the last line of this .tex file.

\begin{document}

% The following commands are self explanatory. Insert your title, author and date into each command's curly bracket. You can include the abstract, paper header and paper footer information in this section then conclude the section with the command \maketitle (shown as the last line at the end of this section).

\title{\LaTeX\ Tutorial}
\author{Shanaya Malik}

\date{}		% leaving the brackets empty omits the date
% To input the current date, you can type: \date{\today}

% The paper headers
\markboth{\LaTeX\ Tutorial}%
{Moulds \MakeLowercase{\textit{et al.}}: CMPE185}
% The only time the second header will appear is for the odd numbered pages
% after the title page when using the twoside option.
% 
% *** Note that you probably will NOT want to include the author's ***
% *** name in the headers of peer review papers.                   ***
% You can use \ifCLASSOPTIONpeerreview for conditional compilation here if
% you desire.

% The publisher's ID mark at the bottom of the page is less important with
% Computer Society journal papers as those publications place the marks
% outside of the main text columns and, therefore, unlike regular IEEE
% journals, the available text space is not reduced by their presence.
% If you want to put a publisher's ID mark on the page you can do it like
% this:
\IEEEpubid{0000--0000/00\$00.00~\copyright~2007 IEEE}
% or like this to get the Computer Society new two part style.
%\IEEEpubid{\makebox[\columnwidth]{\hfill 0000--0000/00/\$00.00~\copyright~2007 IEEE}%
%\hspace{\columnsep}\makebox[\columnwidth]{Published by the IEEE Computer Society\hfill}}
% Remember, if you use this you must call \IEEEpubidadjcol in the second
% column for its text to clear the IEEEpubid mark (Computer Society jorunal
% papers don't need this extra clearance.)


% use for special paper notices
%\IEEEspecialpapernotice{(Invited Paper)}

\IEEEcompsoctitleabstractindextext{%
\begin{abstract}
%\boldmath
The tutorial offers an introduction to \LaTeX, a typesetting system that is particularly relevant within the academic community and often required for technical documentation. The key topics include handling of text, figures, tables, mathematical equations, and references.  The tutorial also emphasizes the importance of \LaTeX in ensuring document portability and the consistent representation of complex formats across various platforms.
\end{abstract}

% Note that keywords are not normally used for peerreview papers.
\begin{IEEEkeywords}
CSE185, \LaTeX\ Tutorial, IEEEtran, journal, \LaTeX, paper, template.
\end{IEEEkeywords}}

\maketitle

%----- The SECTION Environment -------------------------------------------------------------------

\section{Introduction}

\IEEEPARstart{T}{his} tutorial is designed to introduce novice engineering students to \LaTeX\, a software system for typesetting documents widely used in academia and the professional engineering world. The tool operates on the principle that the writer focuses primarily on the content, using plain text and commands to control formatting.

\subsection{Why is this Tutorial Helpful?}
\LaTeX\ is meant for managing complex mathematical equations, tables, and figures, making it an ideal choice for technical documents. The tutorial will guide students through the process of creating documents that are visually professional and precisely formatted.

\subsubsection{Why Should I Learn \LaTeX{}?}
If students are familiar with \LaTeX, it will allow for the development of well-structured scientific reports and papers. \LaTeX\ documents maintain consistency across different computing platforms, improving project collaboration. 

%---- Additional Features ----
\section{Creating a \LaTeX\ File}
\IEEEPARstart{I}{n} order to create a '.tex' file, it is required to use a text editor, but no specialized software is needed. Editors, such as Overleaf, include \LaTeX\ syntax highlighting and integration.

\subsection{Environments}
Environments in \LaTeX\ determine text formatting, and are represented as \texttt{\textbackslash begin} and \texttt{\textbackslash end} commands. It can be used for various purposes, such as itemizing lists, aligning text, or formatting mathematical expressions.

\subsubsection{Begin and End}
The \texttt{\textbackslash begin\{environment\}} command initiates the environment, where \texttt{environment} is the name of the environment you wish to use.  The \texttt{\textbackslash end\{environment\}} command signals the end of an environment. \LaTeX\ will stop applying the special formatting of the environment after this command.  

\subsection{Reserved Characters}
In \LaTeX\, there are some characters have special meanings and are 'reserved' for specific commands and functions. The characters include symbols, such as the backslash `\`, which indicates the start of a command, and the percent `\%`, which comments out text, to name a few. In order to display these characters as part of your text, the character must be escaped using specific commands. 

\subsubsection{Why are these Reserved?}
There are reserved characters in \LaTeX\, which serve specific functions in the typesetting process, such as invoking commands or formatting text. The characters cannot be used directly to display the literal symbols without being escaped.

\subsubsection{Explain the Functions of the Characters: \textbackslash , \textasciitilde , \textbackslash\textbackslash , \%}

\begin{itemize}
    \item \textbackslash: Introduces a command or a \LaTeX keyword.
    \item \textasciitilde: Produces a non-breakable space.
    \item \textbackslash\textbackslash: Creates a new line.
    \item \% : Starts a comment.
\end{itemize}

\subsubsection{What if you Want to Display these Characters?}
In order to display reserved characters in the text, a specific set of commands tells \LaTeX to treat those as regular characters and not as commands or special symbols.

\begin{itemize}
    \item To display a backslash, use \texttt{\textbackslash textbackslash}.
    \item To display a tilde, use \texttt{\textbackslash textasciitilde} or place it inside math mode as \texttt{\$ \textasciitilde \$}.
    \item To display double backslashes, use \texttt{\textbackslash textbackslash\textbackslash textbackslash}.
    \item To display the percent sign, use \texttt{\textbackslash textpercent}.
\end{itemize}

\subsection{Preamble}
The preamble of a \LaTeX document is the section where you define the global settings and include the necessary packages that affect the entire document. It starts right after the \textbackslash documentclass command and continues until the \textbackslash begin\{document\} command. The preamble sets up elements such as the document type, packages, custom commands, and other configurations needed.

\subsubsection{\textbackslash documentclass[] and class files}
The \textbackslash documentclass[] command is the first line in a \LaTeX\ file, and it defines the layout and style of the document. The options within the square brackets set global parameters such as font size, paper size, and whether the document is one-sided or two-sided. There are also class files, which are predefined settings that format the document according to the type of content. 

\subsubsection{Packages}

Packages in \LaTeX\ are extensions that add functionality to the typesetting system. The packages are included in the preamble using the \textbackslash usepackage\{name\} command and provide additional fonts and macros or enable the document to handle complex tasks, such as advanced mathematical formatting, inserting images, handling links, or modifying the layout. 

\subsubsection{Begin and End Document}
The \textbackslash begin\{document\} command marks the beginning of the document's content section. If there is any text before this command, it is considered to be a part of the preamble.  In comparison, the \textbackslash end\{document\} command is used to signify the end of the document. If there is anything written after this command, it will not be processed.

\subsection{Title and Heading Information}
The title, author, and date are elements declared in the preamble and displayed in the document through a specific command.

\subsubsection{\textbackslash maketitle command}
If the title, author, and date have been set in the preamble, it can use the `\maketitle` command within the document body to generate a title page with this information. The `\textbackslash{}maketitle` command uses the information provided by `\textbackslash{}title`, `\textbackslash{}author`, and `\textbackslash{}date` to format a standard title page according to the document class being used. The `\textbackslash{}maketitle` command is flexible and will automatically adjust the formatting of the title page based on the document class and any packages that might influence title page formatting. 

\section{Sections}
\IEEEPARstart{I}{n} a \LaTeX\ document, sections are major divisions and are created using the \textbackslash{}section command. This command numbers the sections and includes it in the document’s Table of Contents.

\subsection{Subsections}
There are subsections, which are used to create smaller divisions within a section. In \LaTeX\, subsections are created using the \textbackslash{}subsection command.

\section{Body Text: Paragraphs and Content}
\IEEEPARstart{T}{he} body text of a document is essential as it includes the main message(s) and information. In a \LaTeX text document, paragraphs are fundamental, and created by leaving a blank line between two blocks of text. The behavior can be altered by changing the length of \texttt{\textbackslash parindent} for indentation and \texttt{\textbackslash parskip} for vertical spacing between paragraphs.  

\subsubsection{Indentation and Spacing}
The text in \LaTeX\ is justified as the default, meaning that the lines stretch to be left-aligned (\texttt{l}), centered (\texttt{c}), or right-aligned (\texttt{r}).

\subsection{Content Blocks and Quotations}
The \texttt{quote} and \texttt{quotation} environments can be used. These provide a visual distinction from the rest of the text, indenting the content from the main body.

\section{Tables}
\IEEEPARstart{T}{he} creation and management of tables in \LaTeX\ involve a combination of the \texttt{table} and \texttt{tabular} environments.

\subsection{Difference Between Table and Tabular Environments}
The \texttt{\textbackslash begin\{table\}} environment is used to add a floating table structure, which can include a caption or a label for referencing. This is particularly useful when the document's layout and formatting require the table not to disrupt the text. The \texttt{\textbackslash begin\{tabular\}} environment is where the table data is inputted and formatted. 

\subsection{Characters and Symbols in Tables}
In the \texttt{tabular} environment, the following symbols and characters are often used:
\begin{itemize}
    \item \texttt{l}, \texttt{c}, \texttt{r} represent column text alignment (left, center, right).
    \item \texttt{|} represents a vertical line between columns.
    \item \texttt{\textbackslash hline} is used to insert a horizontal line between rows.
    \item \texttt{\textbackslash\textbackslash} ends a row and starts a new one.
\end{itemize}

\subsection{Entering Content}
In order to enter content into a table, each cell is separated by an ampersand (\texttt{\&}), and each row is concluded with a double backslash (\texttt{\textbackslash\textbackslash}). 

\subsection{Process of Creating a Table}
The following is an example of a basic table in \LaTeX\, which includes two columns and three rows. The \texttt{table} and \texttt{tabular} environments can be used.  

\begin{enumerate}
    \item The \texttt{table} environment should be initiated with a \texttt{\textbackslash begin\{table\}} command, to create a floating table.
    \item The \texttt{\textbackslash centering} command will center the table. 
    \item If the writer would like to add a caption, it can be done with \texttt{\textbackslash caption\{...\}}, whereas a label for referencing will start with \texttt{\textbackslash label\{...\}}.
    \item The \texttt{tabular} environment can begin with \texttt{\textbackslash begin\{tabular\}\{...\}}, defining the alignment for each column.  
    \item In order to insert horizontal lines before and after the header row, use \texttt{\textbackslash hline}. 
    \item Enter the header row, bold the text with \texttt{\textbackslash textbf}, and separate columns with \texttt{\&}.
    \item End each row with \texttt{\textbackslash\textbackslash} and add \texttt{\textbackslash hline} to insert a horizontal line after the row.
    \item The \texttt{table} and \texttt{tabular} environments can be closed with \texttt{\textbackslash end\{table\}} and \texttt{\textbackslash end\{table\}} respectively.
\end{enumerate}

\begin{table}[h]
\centering
\caption{An Example Table}
\label{tab:example}
\begin{tabular}{|l|c|}
\hline
\textbf{Left-Aligned Header} & \textbf{Centered Header} \\ \hline
Data 1                       & Analysis 1            \\ \hline
Data 2                       & Analysis 2            \\ \hline
Data 3                       & Analysis 3            \\ \hline
\end{tabular}
\end{table}

\section{Figures}
\IEEEPARstart{F}{igures} are an important aspect of scientific documentation, as it allows for the visual representation of data.  

\subsection{Including the Graphicx Package}
The \texttt{graphicx} package is relevant for managing images in \LaTeX. It should be included in the preamble, enabling commands such as \texttt{\textbackslash includegraphics}. The following line can be added at the beginning of the \LaTeX document, before \texttt{\textbackslash begin\{document\}}:
\begin{verbatim}
\usepackage{graphicx}
\end{verbatim}

\subsection{Inserting a Figure}
In order to add a figure to your document, use the \texttt{figure} environment. The environment allows the writer to include images, provide captions, and ensure figures are correctly positioned.

\begin{itemize}
  \item The figure environment with placement specifiers is started with \texttt{\textbackslash begin\{figure\}[ht]}. 
  \item The \texttt{\textbackslash centering} command centers the figure in the text column.
  \item In order to insert the image, the command, \texttt{\textbackslash includegraphics\{image.png\}}, should be used, where the path is replaced with the image file path. 
  \item The following command will allow the writer to insert a caption, \texttt{\textbackslash caption\{...\}}.
  \item In the text, the figure can be referenced with \texttt{\textbackslash label\{...\}}.
\end{itemize}

\subsubsection{Inserting a Titration Plot}
\begin{figure}[htbp]
  \centering
  \includegraphics[width=\linewidth]{Titration_Curve.png}
  \caption{Table of pH Measurements at Incremental Volumes of Base Added During a Titration Process}
  \label{fig:titrationplot}
\end{figure}

\subsection{Creating Graphs}
In \LaTeX, graphs can be created using the \texttt{pgfplots} package, or can include a pre-made graph using \texttt{\textbackslash includegraphics}.

\subsection{Including a Meaningful Caption}
The usage of captions is that it provides context and explain the significance of the figure or graph, which can be initiated with \texttt{\textbackslash caption} and followed with the description.

\subsection{Referencing the Graph in Text}
In order to reference the figure in the text, it can be done using the \texttt{\textbackslash ref} command with the label assigned to the figure.

\section{Mathematical Formulas}

\IEEEPARstart{M}{athematical} formulas are central to engineering and scientific documents. \LaTeX provides a system for typesetting complex mathematical equations and symbols.  

\subsection{Equation Environments}
\LaTeX uses two primary methods for displaying equations: inline and display. Inline equations appear within the text, while display equations are separated from the main text.

\subsubsection{Inline Equations}
Insert inline equations using dollar signs (\$). Einstein's energy-mass equivalence formula can be written inline like the following. 

\( E = mc^2 \).

\subsubsection{Display Equations}
The display equations are separated from the text and often centered. The `equation` environment is used for that purpose.

\begin{verbatim}
\begin{equation}
    E = mc^2
\end{equation}
\end{verbatim}

It will appear in the \LaTeX document like this:

\begin{equation}
    E = mc^2
\end{equation}

\subsection{Symbols}
\LaTeX supports an extensive array of mathematical symbols, including Greek letters \( \delta \); summation symbols \( \sum \); and integral symbols \( \int \).

\subsection{Fractions}
Fractions in LaTeX are created with the \textbackslash frac{} command, which takes the numerator and the denominator as arguments. 

\begin{verbatim}
\frac{numerator}{denominator}
\end{verbatim}

\subsubsection{Writing the Fraction One-Half}

In order to write one-half as a fraction, the following command should be used in the \LaTeX document. 
\begin{verbatim}
\frac{1}{2}
\end{verbatim}

This will appear as \( \frac{1}{2} \). 

\subsection{Complex Equations}
\LaTeX provides environments for complex equations involving matrices or multiple aligned equations.

\subsubsection{Matrices}
A 2x2 matrix would be written as the following. 

\begin{verbatim}
\begin{equation}
\begin{matrix}
a & b \\
c & d \\
\end{matrix}
\end{equation}
\end{verbatim}

\subsubsection{Aligned Equations}
For equations that should be aligned at the equal sign, it would be written like this. 

\begin{verbatim}
\begin{align}
y &= 2x + 3 \\
y &= -x - 6
\end{align}
\end{verbatim}

\section{How to: Acknowledgments}

\IEEEPARstart{T}{he} acknowledgments section is an important part of any academic paper or report. It is a space where the author can thank those who have contributed to the work but are not among the authors, including advisors; funding agencies; technical staff; or colleagues, to name a few. 

In \LaTeX, an acknowledgments section can be included through using the \texttt{\textbackslash section\{Acknowledgments\}} command. The section often comes just before the references section.

The use of the \texttt{\textbackslash section*\{Acknowledgments\}} command allows the writer to create a non-numbered section. The asterisk (*) next to \texttt{\textbackslash section} prevents the section number from appearing next to the title.

\subsection*{Placement of Acknowledgments}

In a document, the acknowledgments should ideally appear after the conclusion and prior to the references. The placement ensures that the acknowledgments are easily found by readers but do not disrupt the flow of the paper's main content.

\section{How to: References}

\IEEEPARstart{M}{anaging} references in \LaTeX{} is efficient and allows for consistent citation formatting. This guide covers the \texttt{thebibliography} environment and key commands: \texttt{\textbackslash label}, \texttt{\textbackslash ref}, and \texttt{\textbackslash cite}.

\subsection{Thebibliography Environment}
Listing references can be done within the the \texttt{thebibliography} environment, where each entry is defined by \texttt{\textbackslash bibitem}. The \texttt{thebibliography} environment acts as a list, where each reference is an item.

\subsection{Using \textbackslash label and \textbackslash ref}
In order to assign a label to figures, tables, or sections, please use the \texttt{\textbackslash label} command, and reference it using \texttt{\textbackslash ref}.

\subsection{Citing References}
Use \texttt{\textbackslash cite} for citations in your text, referencing the label of the \texttt{\textbackslash bibitem}. 

\subsection{Compiling the Document}
In order to process references, the document should be compiled twice. The first pass collects the reference keys, and the second pass updates the citations.

\section{Conclusion}

\IEEEPARstart{T}{he} use and flexibility of \LaTeX allows for an improved form of technical documentation. The capabilities of the typesetting program allow for any engineer to become better prepared for both academic and professional communication. \LaTeX is an important part of being able to draft technical reports, theses, and research papers. 

\section*{Acknowledgments}

\IEEEPARstart{T}{he} author would like to express gratitude to the individuals and resources that were helpful in the completion of this assignment. This includes the teaching assistant, Kevin Bowden, of the CSE 185S class at UCSC, whose guidance has been valuable in better introducing the topic of the assignment. Professor Gerald Moulds has also provided insightful instruction in the Technical Writing for Engineers course, which has improved the author's understanding and application of technical communication principles, particularly in the subject. In addition, Professor David Draper of the STAT 131 class, Introduction to Probability Theories, at UCSC, provided a brief introduction to \LaTeX\ allowed the author to discover the tool for scientific documentation. The creators of Overleaf have provided a comprehensive tutorial that, while not directly cited, was read and appreciated for the clear explanation and practical examples that facilitated a deeper understanding of how to use \LaTeX\ for document preparation.

\begin{thebibliography}{1}

\bibitem{OverleafTutorials}
Overleaf. (n.d.). \textit{Learn LaTeX in 30 minutes}. Retrieved from \url{https://www.overleaf.com/learn/latex/Tutorials}

\end{thebibliography}


\end{document}

%----- Additional Features -----------------------------------------------------------------------

\section{Additional Features}

\subsection{Figures}
% FIGURES:

% Note the FIGURE Environment created by the \begin{figure} and \end{figure} commands.

\begin{figure}[h] 	% There are several different modifiers that can be used in [].
\centering
\includegraphics[width=1.5in]{slug.pdf}
\caption{Sammy the Slug}
\label{fig_slug}
\end{figure}

% You will need to use appropriate file types for figures and will also need to include that image file in the same folder as your .tex file. 

%-------------------------------------------------------------------------------------------------
\subsection{Label, Cite, and Ref Commands}
You will need to be able to define and explain how to use the following commands:
\begin{verbatim} \label{fig_slug} \end{verbatim} as it appears in the figure environment\\
\begin{verbatim} \cite{IEEEhowto:kopka} \end{verbatim} appears like: \cite{IEEEhowto:kopka}\\
\begin{verbatim} \ref{fig_slug} \end{verbatim} appears like: \ref{fig_slug}\\

% IMPORTANT NOTE: In order to assign the correct reference number to each label, you may have to compile your code twice. 

%-------------------------------------------------------------------------------------------------
\subsection{Tables}
An example of a floating table.

% An example of a floating table. Note that, for IEEE style tables, the 
% \caption command should come BEFORE the table. Table text will default to
% \footnotesize as IEEE normally uses this smaller font for tables.
% The \label must come after \caption as always.
%
\begin{table}[h]
%% increase table row spacing, adjust to taste
\renewcommand{\arraystretch}{1.3}
% if using array.sty, it might be a good idea to tweak the value of
%\extrarowheight as needed to properly center the text within the cells
\caption{An Example of a Table}
\label{table_example}
\centering
%% Some packages, such as MDW tools, offer better commands for making tables
%% than the plain LaTeX2e tabular which is used here.
\begin{tabular}{|c||c|}
\hline
One & Two\\
\hline
Three & Four\\
\hline
\end{tabular}
\end{table}


% Note that IEEE does not put floats in the very first column - or typically
% anywhere on the first page for that matter. Also, in-text middle ("here")
% positioning is not used. Most IEEE journals use top floats exclusively.
% However, Computer Society journals sometimes do use bottom floats - bear
% this in mind when choosing appropriate optional arguments for the
% figure/table environments.
% Note that, LaTeX2e, unlike IEEE journals, places footnotes above bottom
% floats. This can be corrected via the \fnbelowfloat command of the
% stfloats package.


\section{Conclusion}
Conclusion goes here.

%----- APPENDICES --------------------------------------------------------------------------------
\appendices
\section{Appendix Title}
Appendix one text goes here.

% you can choose not to have a title for an appendix
% if you want by leaving the argument blank
\section{}
Appendix two text goes here.


%----- ACKNOWLEDGEMENT SECTION -------------------------------------------------------------------
% Explain what the asterisk * does in the next line: 
\section*{Acknowledgements}

The author would like to thank...\\ \\

Reminder: you will need to explain how to include an Acknowledgement Section and then include your own Acknowledgement Section at the end of your own tutorial. Same applies for the References/Bibliography.


%----- BIBLIOGRAPHY ------------------------------------------------------------------------------

% You will need to explain how to include the bibliography section as follows. Explain the environment and how to add new items.
% Including how \ref, \cite and \label should be included here.

% Reminder: you will need to explain how to include the Bibliography Section and then include your own Bibliography at the end of your own tutorial.

\begin{thebibliography}{1}

\bibitem{IEEEhowto:kopka}
H.~Kopka and P.~W. Daly, \emph{A Guide to {\LaTeX}}, 3rd~ed.\hskip 1em plus
  0.5em minus 0.4em\relax Harlow, England: Addison-Wesley, 1999.

\end{thebibliography}

%----- Optional: BIOGRAPHY Section ---------------------------------------------------------------
 
% If you have an EPS/PDF photo (graphicx package needed) extra braces are
% needed around the contents of the optional argument to biography to prevent
% the LaTeX parser from getting confused when it sees the complicated
% \includegraphics command within an optional argument. (You could create
% your own custom macro containing the \includegraphics command to make things
% simpler here.)
%\begin{biography}[{\includegraphics[width=1in,height=1.25in,clip,keepaspectratio]{mshell}}]{Gerald Moulds}
% or if you just want to reserve a space for a photo:

\begin{IEEEbiography}{Gerald Moulds}
Biography text here.
\end{IEEEbiography}

% if you will not have a photo at all:
\begin{IEEEbiographynophoto}{John Doe}
Biography text here.
\end{IEEEbiographynophoto}

% insert where needed to balance the two columns on the last page with
% biographies
%\newpage

\begin{IEEEbiographynophoto}{Jane Doe}
Biography text here.
\end{IEEEbiographynophoto}

% You can push biographies down or up by placing
% a \vfill before or after them. The appropriate
% use of \vfill depends on what kind of text is
% on the last page and whether or not the columns
% are being equalized.

%\vfill

% Can be used to pull up biographies so that the bottom of the last one
% is flush with the other column.
%\enlargethispage{-5in}

\end{document}
